% 08-invariant — что именно сохраняет дискретный закон (инвариант, симметрия, предел M/T)
\chapter{Что именно сохраняется: биекция, \texorpdfstring{$\mathbb{Z}_4$}{Z4} и предел \texorpdfstring{$U(1)$}{U(1)}}
\label{ch:invariant}

Всякая физика держится на законах сохранения, а законы сохранения --- на симметриях.
Вопрос, который обычно не задают: \emph{что} именно сохраняет дискретный закон,
если у него нет ни непрерывного времени, ни метрики, ни нормы волновой функции.
Эта глава --- ответ, и он не совпадает с тем, что записано в~\Axref{ax:A3}.
Ниже показано, из чего следует сохранение информации на самом деле,
какая группа симметрии у шага точная, и где проходит граница между~$\Mlayer$ и~$\Tlayer$.

\section{Как запись появилась: глобальная норма до спуска}

Условие~\Axref{ax:A3} было введено в форме, привычной по квантовой механике:
норма $N(t)=\sum_x z^\dagger z$ есть инвариант эволюции, линейная часть закона ---
произведение локальных унитарных операторов, нелинейная --- умножение на фазу.
Источник требования физический: в квантовой механике эволюция унитарна,
$\langle\wavePsi|\wavePsi\rangle=\mathrm{const}$.

Существенно, \emph{когда} это писалось. Тогда квантовая механика ещё бралась
в глобальной форме: унитарный оператор на гильбертовом пространстве,
антисимметрия волновой функции как постулат о целом, коллапс как нечто нелокальное.
Всё это --- утверждения про целое, и норма в них --- интеграл по всему пространству.

Спуск (гл.~\ref{ch:descent}) и вычислительная лаборатория серии поменяли картину.
Программа «квантовая механика из локального детерминированного автомата» поставлена
не здесь: она принадлежит Г.~’т~Хоофту (\cref{sec:thooft-programme}),
и там же сказано, в чём он ошибся --- в решётке.
Пять его открытых проблем разобраны в томе «Вычисления», в разделе о них.
Там же доказано то, на что опирается эта глава: шаг $g$ биективен
(том «Вычисления», теорема о побитовой обратимости модулярной «лягушки»), а неинъективная склейка микросостояний
в наблюдаемые шаблоны образует четырёхуровневую иерархию классов информационной
эквивалентности (таблица классов того же тома) ---
это и есть точный ответ на вопрос, что стоит на месте нормы.
Новое здесь --- не принцип локальности, а то, что он доведён до конкретных ворот
с числами: запрет Паули как локальные ворота в одной ячейке~$\PlanckCell$
($\Phi_{\mathrm{Pauli}}=256$ тик, полуоборот кольца), зеркало
$Z_{\mathrm{next}}=-Z_{\mathrm{past}}$ из сокращения члена~$2Z$,
локальная мёртвая зона Гейзенберга и световой конус как источник видимой нелокальности.

Отсюда первый вывод главы: \textbf{глобальная норма есть тень локального правила,
а не наоборот.} Условие~\Axref{ax:A3} в записанной форме относится к тому языку,
с которого начался спуск, и на носителе ему негде жить: на кольце
$\mathbb{Z}_{\phaseRingCount}$ величины $|z|^2$ как числа не существует,
она зависит от выбора представителя.

\section{Что сохраняется на самом деле: биекция}

\begin{theorem}[Сохранение информации без метрики]
\label{thm:info-bijection}
Пусть $\mathcal{C}$ --- пространство конфигураций на $\carrierLattice$
с конечным алфавитом $\mathbb{Z}_{\phaseRingCount}[i]^2$ на узле,
а $g:\mathcal{C}\to\mathcal{C}$ --- закон эволюции, удовлетворяющий~\Axref{ax:A13}.
Тогда:
\begin{enumerate}[label=(\roman*)]
\item никакая информация о начальном состоянии не теряется: $g$ --- биекция,
      и $\symCapitalPsi^{t}$ однозначно восстанавливается из $\symCapitalPsi^{t+1}$;
\item траектория строго периодична: период не превосходит $|\mathcal{C}|$,
      и возврат Пуанкаре существует для любой начальной конфигурации;
\item ни одна величина вида $\sum_x f\bigl(z(x)\bigr)$ не обязана сохраняться:
      у биекции на конечном множестве нет нетривиальных первых интегралов,
      кроме метки орбиты.
\end{enumerate}
\end{theorem}

\begin{proof}
\emph{(i)} --- содержание~\Axref{ax:A13}: существует $g^{-1}$.

\emph{(ii)} --- \pfrom{биективности~$g$ на конечном $\mathcal{C}$}
орбита есть подмножество конечного множества, и повторение состояния
влечёт периодичность; оценка периода следует из мощности.

\emph{(iii)} --- \pusing{контрпримером}
достаточно указать биекцию, не сохраняющую $\sum|z|^2$: в $\mathbb{Z}_{\phaseRingCount}[i]$
умножение на неединичный элемент меняет модуль, а сложение с ненулевым сдвигом
меняет сумму. Общего механизма, который заставлял бы такую величину сохраняться,
у биекции нет.
\end{proof}

\begin{corollary}[Теорема~\ref{thm:main-23} не нуждается в норме]
\label{cor:heat-death-without-norm}
Вывод теоремы о невозможности тепловой смерти~$\Mlayer$ (\cref{thm:main-23})
сохраняется целиком: отсутствие останова и отсутствие необратимого распада
следуют из пунктов (i)--(ii) теоремы~\ref{thm:info-bijection}
и нефизичности вырожденной неподвижной точки (\Axref{ax:A5}).
Доказательство меняется: лемма~\ref{lem:232} перевешивается с~\Axref{ax:A3}
на~\Axref{ax:A13}.
\end{corollary}

Это не ослабление, а усиление: из теоремы уходит континуальная величина,
и остаётся свойство конечной перестановки.

\section{Какая группа симметрии у шага точная}

\begin{proposition}[Точная группа симметрии зависит от записи регистра]
\label{prop:z4-exact}
В декартовой записи $Z=U+iV\in\mathbb{Z}_{\phaseRingCount}[i]$ модуль сохраняют только
умножения на единицы $\pm1,\ \pm i$. Для любого другого множителя $\alpha$
существует $z$ с $|\alpha z|^2\ne|z|^2$, поэтому там точная группа --- $\mathbb{Z}_4$,
а произвольный поворот реализуется лишь приближённо.
В полярной записи $(\mu,\varphi)\in\mathbb{Z}_{2^{\mathrm{frac\_bits}}}\times\mathbb{Z}_{\phaseRingCount}$
(внутренние координаты регистра) поворот есть сложение фаз по модулю $\phaseRingCount$
и точен для всех $\phaseRingCount$ значений: точная группа --- $\mathbb{Z}_{\phaseRingCount}$.
\end{proposition}

\begin{proof}
\emph{Декартово.}
\pfrom{мультипликативности модуля в $\mathbb{Z}[i]$}
$|\alpha z|^2=|\alpha|^2|z|^2$; равенство при всех $z$ требует $|\alpha|^2=1$,
то есть $\alpha\in\{\pm1,\pm i\}$.
\pfollows
всякий иной поворот есть аппроксимация, и её точность ограничена разрешением
таблицы поворотов.

\emph{Полярно.}
\pusing{определения внутренних координат регистра}
фаза есть элемент $\mathbb{Z}_{\phaseRingCount}$, а поворот --- сложение
$\varphi\mapsto\varphi+k$ по модулю $\phaseRingCount$; сложение в кольце вычетов
биективно при любом $k\in\mathbb{Z}_{\phaseRingCount}$, и ни одно значение не выделено.
\end{proof}

\begin{remark}[$\mathbb{Z}_4$ --- свойство записи, а не носителя]
\label{rem:u1-is-t-limit}
Отсюда меняется прочтение~\Axref{ax:A4}. На~$\Mlayer$ в полярном регистре
$\mathbb{Z}_{\phaseRingCount}$ все углы различимы и поворот точен, так что дискретная
симметрия носителя не беднее непрерывной: континуальная $U(1)$ есть предел
$\mathbb{Z}_{\phaseRingCount}$ при $\phaseRingCount\to\infty$, а не отдельный объект~$\Tlayer$.
Деградация до $\mathbb{Z}_4$ возникает только при декартовой записи через таблицу
поворотов --- и именно она даёт ложное впечатление, что симметрия носителя урезана.
\end{remark}

\section{Тень инварианта: квадратичный вычет}

\begin{definition}[Кольцевой квадрат]
\label{def:ring-square}
Для дорожек $u,v\in\mathbb{Z}_{2^{b}}$ компоненты спинора положим
\[
  q(z) := (u^2+v^2)\ \mathrm{mod}\ 2^{b}.
\]
\end{definition}

\begin{proposition}[$q$ не зависит от разреза]
\label{prop:q-cut-free}
Сдвиг представителя на полный период, $u\mapsto u+2^{b}$, не меняет $q$:
$2\cdot2^{b}u+(2^{b})^2\equiv0\pmod{2^{b}}$.
Следовательно $q$ есть функция класса, а не координаты, и это единственная
величина, которая ведёт себя как модуль, но определена на кольце честно.
\end{proposition}

\begin{proposition}[$q$ инвариантен для поворота и не инвариантен для шага]
\label{prop:q-scope}
Точный поворот (умножение на единицу) сохраняет $q$ поячеечно.
Шаг эволюции содержит, кроме поворота, сложение по схеме «лягушки»,
поэтому $q$ поячеечно не сохраняется: измерение на сиде с открытыми воротами
даёт сохранение $q$ в $94\%$ ячеек за $32$ такта.
\end{proposition}

Таким образом $q$ --- не закон сохранения, а \emph{тень} симметрии:
там, где шаг сводится к повороту, она точна; где нет --- дрейфует с шагом,
задаваемым разрешением таблицы.

\section{Какая форма закона сохранения следует: расщепление голономии и закон движения \texorpdfstring{$\mu$}{mu}}
\label{sec:holonomy-split-mu}

Из локального закона следует не глобальная константа, а \textbf{уравнение непрерывности}:
изменение веса ячейки $q=\mu^2$ равно потоку через её двенадцать связей.
Чтобы увидеть, откуда берётся этот поток и какую роль выполняет амплитуда
$\mu\in\{1,\dots,2^{\mathrm{frac\_bits}}\}$ рядом с фазой $\varphi_{\mathrm{disc}}\in\mathbb{Z}_{\phaseRingCount}$,
достаточно записать действие дискретного лапласиана звезды $N(x)$ в полярной форме $z=\mu e^{i\varphi}$
и умножить уравнение эволюции $i\,\partial_t z = -\gamma\,\Delta_{\mathrm{disc}} z$ на $z^*(x)$.

\begin{theorem}[Ортогональное расщепление голономии и точное уравнение непрерывности]
\label{thm:holonomy-split-continuity}
Для любого поля $z(x)=\mu(x)e^{i\varphi(x)}$ на решётке с симметричным шаблоном соседей
$y\in N(x)$ ($y\in N(x)\iff x\in N(y)$) произведение $i(\partial_t z)z^*$ и голономия звезды
$\zeta(x)=\sum_{y\in N(x)} z(y)z^*(x) = \zeta_{\mathrm{real}}(x)+i\,\zeta_{\mathrm{imag}}(x)$
расщепляются на действительную (чётную по $\Delta\varphi$) и мнимую (нечётную по $\Delta\varphi$) части:
\begin{equation}
\label{eq:holonomy-split}
  i\,(\partial_t z)\,z^*
  \;=\;
  \underbrace{-\mu^2\,\partial_t\varphi}_{\text{действительная часть}}
  \;+\;
  i\,\underbrace{\frac{1}{2}\,\partial_t(\mu^2)}_{\text{мнимая часть}}
  \;=\;
  -\gamma\Bigl(\bigl(\zeta_{\mathrm{real}} - |N|\,\mu^2\bigr) + i\,\zeta_{\mathrm{imag}}\Bigr).
\end{equation}
Отсюда следуют два ортогональных закона --- для фазы $\varphi$ и для амплитуды $\mu$:
\begin{enumerate}[label=(\roman*)]
\item \textbf{Чётный дефицит голономии задаёт положительную угловую частоту фазы $\omega(x) = -\partial_t\varphi(x)$ при фиксированном $\mu$:}
  \begin{equation}
  \label{eq:even-holonomy-phase}
  \begin{aligned}
    -\mu^2\,\partial_t\varphi
    &\;=\;
    \mu^2\,\omega(x)
    \;=\;
    \gamma\,\delta\zeta_{\mathrm{real}}(x), \\
    \delta\zeta_{\mathrm{real}}(x)
    &\;:=\;
    |N|\,\mu^2(x) - \zeta_{\mathrm{real}}(x)
    \;=\;
    \sum_{y\in N(x)} \mu(x)\mu(y)\bigl(1 - \cos(\varphi(y)-\varphi(x))\bigr) \;\ge\; 0.
  \end{aligned}
  \end{equation}
  Слагаемые противоположных связей $+\mathbf{e}$ и $-\mathbf{e}$ в $\delta\zeta_{\mathrm{real}}$
  имеют одинаковый знак ($1-\cos(+\Delta\varphi)=1-\cos(-\Delta\varphi)\ge 0$) и не сокращаются
  на симметричной звезде: именно $\delta\zeta_{\mathrm{real}}$ задаёт частоту плоской волны,
  равномерный ход фазы кипящего океана и одевающий ореол вихря на первой координационной сфере $S_1$.
\item \textbf{Нечётная голономия $\zeta_{\mathrm{imag}}$ есть дивергенция антисимметричного реберного потока и переносит вес $q=\mu^2$ между соседями:}
  \begin{equation}
  \label{eq:odd-holonomy-continuity}
    \frac{1}{2}\,\partial_t\bigl(\mu^2(x)\bigr)
    \;=\;
    -\gamma\,\zeta_{\mathrm{imag}}(x)
    \;=\;
    -\sum_{y\in N(x)} J(x\to y)
    \;=\;
    -\operatorname{div}_{\mathrm{disc}}\mathbf{j}(x),
  \end{equation}
  где поток по направленному ребру $(x\to y)$ равен
  \begin{equation}
  \label{eq:edge-flux}
    J(x\to y) \;:=\; \gamma\,\mu(x)\mu(y)\sin\bigl(\varphi(x)-\varphi(y)\bigr) \;=\; -J(y\to x).
  \end{equation}
\item \textbf{Источник тождественно равен нулю:} в силу антисимметрии $J(x\to y)=-J(y\to x)$
  на каждом ребре решётки сумма дивергенций по всем узлам равна нулю для любой конфигурации:
  \begin{equation}
  \label{eq:zero-global-source}
    \sum_{x\in\carrierLattice} \zeta_{\mathrm{imag}}(x)
    \;=\;
    -\frac{1}{\gamma}\sum_{x\in\carrierLattice}\operatorname{div}_{\mathrm{disc}}\mathbf{j}(x)
    \;\equiv\; 0.
  \end{equation}
\end{enumerate}
\end{theorem}

\begin{proof}
Подстановка $z=\mu e^{i\varphi}$ даёт $\partial_t z = (\partial_t\mu + i\mu\partial_t\varphi)e^{i\varphi}$,
откуда $i(\partial_t z)z^* = -\mu^2\partial_t\varphi + i\mu\partial_t\mu = -\mu^2\partial_t\varphi + \frac{i}{2}\partial_t(\mu^2)$.
Для звезды соседей $z(y)z^*(x) = \mu(x)\mu(y)e^{i(\varphi(y)-\varphi(x))}$, поэтому
$\zeta_{\mathrm{real}}(x)-|N|\mu^2(x) = \sum_{y\in N(x)}\bigl(\mu(x)\mu(y)\cos(\varphi(y)-\varphi(x))-\mu^2(x)\bigr)$
и $\zeta_{\mathrm{imag}}(x) = \sum_{y\in N(x)}\mu(x)\mu(y)\sin(\varphi(y)-\varphi(x))$.
Приравнивание действительных и мнимых частей даёт~\eqref{eq:even-holonomy-phase}--\eqref{eq:edge-flux}.
Суммирование~\eqref{eq:edge-flux} по всем ориентированным рёбрам $(x,y)$ решётки разбивается
на пары противоположных ориентаций $J(x\to y)+J(y\to x)=0$, что доказывает~\eqref{eq:zero-global-source}.
\end{proof}

\begin{corollary}[Что делает амплитуда $\mu$ на $\Mlayer$]
\label{cor:what-mu-does}
В полярном регистре $(\mu,\varphi_{\mathrm{disc}})\in\{1,\dots,64\}\times\mathbb{Z}_{512}$ амплитуда $\mu$
выполняет три согласованные функции:
\begin{enumerate}[label=(\arabic*)]
\item \textbf{Инерция фазы (почему ядро возбуждается выше порога Гейзенберга, а массивная оболочка преломляет внешние волны):}
  приращение фазы $\Phi$ определяется отношением $\delta\zeta_{\mathrm{real}}(x)/D(x)$, где
  в числителе стоит фазовое натяжение $\delta\zeta_{\mathrm{real}}=|N|\mu^2-\zeta_{\mathrm{real}}$,
  а в знаменателе --- полный вес ворот $D(x)=\zeta_{\mathrm{real}}(x)+\mu^2(x)+K_P$.
  Для \emph{собственной} фазы вихря на ядре и первой сфере $S_1$ противоположные соседи находятся
  в противофазе ($\cos\Delta\varphi<0$), поэтому дефицит $\delta\zeta_{\mathrm{real}}$ максимален
  ($18648\ \text{квантов}^2$ в ядре и $8417\ \text{квантов}^2$ на $S_1$ против $8\ \text{квантов}^2$ в далёком вакууме)
  и открывает затвор Гейзенберга ($|\Phi|\ge 41$ тик, $n_E\ge 1$).
  Напротив, для \emph{внешней пробной волны} с малым градиентом $\delta\varphi_{\mathrm{test}}$,
  проходящей сквозь одетую оболочку массивного дефекта, знаменатель $D(r)=\zeta_{\mathrm{real}}(r)+\mu^2(r)+K_P$
  монотонно растёт при приближении к центру ($D_{\mathrm{vac}}=82 \to \langle D\rangle_{S_3}=2077 \to \langle D\rangle_{S_1}=24138$),
  уменьшая угловую скорость набега фазы пробной волны ($\Delta\Phi_{\mathrm{test}}\propto D(r)^{-1}$).
\item \textbf{Проводимость связи:} и фазовое натяжение~\eqref{eq:even-holonomy-phase},
  и реберный поток~\eqref{eq:edge-flux} пропорциональны произведению $\mu(x)\mu(y)$ на концах связи.
\item \textbf{Перенос веса при столкновениях и стационарность уединённых структур:}
  на стационарных гармонических конфигурациях ($\Delta_{\mathrm{disc}}\varphi=0$: плоская волна,
  линейный рамп кипящего океана, радиально-симметричный уединённый вихрь)
  поток через звезду уравновешен ($\zeta_{\mathrm{imag}}(x)=0$), поэтому профиль амплитуды $\mu(x)$
  стоит на месте, пока фаза вращается со скоростью $\delta\zeta_{\mathrm{real}}(x)>0$.
  При наложении и столкновении волновых пакетов или вихревых пар ($p+m$) гармоничность фазы
  нарушается ($\zeta_{\mathrm{imag}}\ne 0$), и вес $q=\mu^2$ перетекает по рёбрам по закону
  \begin{equation}
  \label{eq:discrete-mu-step}
    q(x,t+1) \;=\; q(x,t) \;-\; \sum_{y\in N(x)} J_{\mathrm{disc}}(x\to y),
    \qquad
    J_{\mathrm{disc}}(y\to x) = -J_{\mathrm{disc}}(x\to y),
  \end{equation}
  между полом~\Axref{ax:A5} ($\mu\ge 1$) и потолком~\Axref{ax:A7} ($\mu\le 64$) без рождения из вакуума.
\end{enumerate}
\end{corollary}

\begin{proposition}[Критерий: поток или источник]
\label{prop:source-or-flux}
Если дискретный шаг меняет вес ячейки по антисимметричному реберному потоку~\eqref{eq:discrete-mu-step},
то $\Delta q(x)+\operatorname{div}_{\mathrm{disc}}\mathbf{j}(x)=0$ поячеечно, и $\sum_x q(x)$ сохраняется точно.
Ненулевой рост $\sum_x q(x)$ возможен только тогда, когда схема подменяет реберный перенос
неунитарным декартовым сложением («лягушкой» на $(U,V)$ без сохранения модуля).
\end{proposition}

\section{Дискретные уравнения Максвелла на комплексе кубооктаэдра и единство взаимодействий без внешних источников}
\label{sec:discrete-maxwell-unification}

Отсутствие источника в~\eqref{eq:zero-global-source} решает вопрос, который в континуальной теории поля
остаётся неразрешимым разделением на «поля-переносчики» и «материю-источник».
В уравнениях Максвелла на~$\Tlayer$ первая пара ($\operatorname{div}\mathbf{B}=0$, $\partial_t\mathbf{B}+\operatorname{rot}\mathbf{E}=0$)
не имеет правых частей, тогда как во второй паре ($\operatorname{div}\mathbf{E}=\rho_{\mathrm{el}}$,
$\operatorname{rot}\mathbf{B}-\partial_t\mathbf{E}=\mathbf{j}_{\mathrm{el}}$) плотность заряда $\rho_{\mathrm{el}}$
и ток $\mathbf{j}_{\mathrm{el}}$ постулируются как внешняя сущность, посаженная в правую часть.
Перевод на клеточный комплекс кубооктаэдра $N_{12}$ ($0$-узлы $x$, $1$-рёбра $e=(x\to y)$,
$2$-грани $f\in 8\triangle\cup 6\square$, $3$-объём $\PlanckCell$) показывает, что на~$\Mlayer$
\textbf{никаких внешних источников нет ни в одном уравнении}: то, что на~$\Tlayer$ читается как «заряд» и «масса»,
есть целочисленный скачок ветви того же самого фазового поля на замкнутых плакетах кубооктаэдра.

\begin{theorem}[Дискретные уравнения Максвелла без внешних источников]
\label{thm:discrete-maxwell-sourceless}
Пусть на $0$-узлах решётки $\carrierLattice$ задан полярный регистр $(\mu(x),\varphi(x))\in\{1,\dots,64\}\times\mathbb{Z}_{\phaseRingCount}$,
и пусть $d_0$ (разность на ребре), $d_1$ (сумма по контуру грани $\partial f$) и $d_2$ (сумма по замкнутой
оболочке $\partial\PlanckCell = 8\triangle\cup 6\square$) --- кограничные операторы клеточного комплекса ($d_{k+1}\circ d_k \equiv 0$).
Определим на направленном ребре $e=(x\to y)$ свёрнутую фазовую связность в главном интервале $(-\phaseRingCount/2,\phaseRingCount/2]$:
\begin{equation}
\label{eq:discrete-edge-connection}
  A_{xy} \;:=\; \mathrm{wrap}\bigl(\varphi(y)-\varphi(x)\bigr)
  \;=\; \bigl(\varphi(y)-\varphi(x)\bigr) - \phaseRingCount\,k_{xy},
  \qquad
  k_{xy}\in\{-1,0,+1\},
\end{equation}
дискретное электрическое поле на ребре $E_{xy} := -\Delta_t A_{xy} = -(d_0\,\Delta_t\varphi)_{xy} + \phaseRingCount\,\Delta_t k_{xy}$
и дискретный магнитный поток через ориентированную грань $f\in 8\triangle\cup 6\square$:
\begin{equation}
\label{eq:discrete-plaquette-flux}
  B_f \;:=\; (d_1 A)_f \;=\; \sum_{e\in\partial f} A_e \;=\; -\phaseRingCount\,(d_1 k)_f \;=\; \phaseRingCount\,n_f,
  \qquad
  n_f \;:=\; -(d_1 k)_f \in \mathbb{Z}.
\end{equation}
Тогда на каждой клетке комплекса выполняются четыре точных дискретных уравнения:
\begin{enumerate}[label=(\roman*)]
\item \textbf{Дискретный закон Фарадея на каждой грани $f\in 8\triangle\cup 6\square$:}
  \begin{equation}
  \label{eq:discrete-faraday}
    \Delta_t B_f + (d_1 E)_f \;=\; \Delta_t(d_1 A)_f - (d_1\,\Delta_t A)_f \;\equiv\; 0.
  \end{equation}
\item \textbf{Дискретный закон Гаусса для магнитного потока на замкнутой оболочке $\partial\PlanckCell = 8\triangle\cup 6\square$:}
  \begin{equation}
  \label{eq:discrete-gauss-b}
    (d_2 B)_{\PlanckCell} \;=\; \sum_{f\in\partial\PlanckCell} B_f \;=\; (d_2 d_1 A)_{\PlanckCell} \;\equiv\; 0,
    \qquad
    \sum_{f\in 8\triangle\cup 6\square} n_f \;\equiv\; 0
    \qquad (d_2\circ d_1 \equiv 0).
  \end{equation}
  Вихревая нить ($n_f\ne 0$) не может оборваться внутри ячейки~$\PlanckCell$: число входящих и выходящих
  квантов потока через $14$ граней кубооктаэдра совпадает тождественно.
\item \textbf{Вторая пара Максвелла как тождество скачков ветви $k_{xy}$:}
  при сопряжённых операторах $d_0^\dagger = -\operatorname{div}_{\mathrm{disc}}$ (сумма по звезде узла $N_{12}(x)$)
  и $d_1^\dagger = \operatorname{curl}_{\mathrm{disc}}$ (сумма по граням, примыкающим к ребру)
  кограница поля раскладывается на гладкую волновую часть и целочисленные дефекты ветви $k_{xy}$ и $n_f = -(d_1 k)_f$:
  \begin{equation}
  \label{eq:discrete-ampere-gauss}
    (d_0^\dagger E)_x \;=\; -(d_0^\dagger d_0\,\Delta_t\varphi)_x + \phaseRingCount\,\Delta_t(d_0^\dagger k)_x,
    \qquad
    (d_1^\dagger B - \Delta_t E)_{xy}
    \;=\;
    \phaseRingCount\,(d_1^\dagger n)_{xy} + (\Delta_t^2 A)_{xy}.
  \end{equation}
  Там, где топологических дефектов нет ($n_f=0$, занятость $b=0$), правая часть равна нулю и по решётке
  распространяются свободные фазовые волны; там, где $n_f=\pm 1$ ($b=1$, планкон), ненулевой ротор
  $\phaseRingCount\,(d_1^\dagger n)_{xy}$ вокруг вихревого ядра и есть то, что на~$\Tlayer$ считывается
  как калибровочный ток и заряд источника.
\item \textbf{Тождественное сохранение заряда без внешних источников:}
  поскольку $d_0^\dagger\circ d_1^\dagger = (d_1\circ d_0)^\dagger \equiv 0$, дивергенция вихревого тока
  тождественно равна нулю на любой конфигурации:
  \begin{equation}
  \label{eq:discrete-charge-conservation}
    d_0^\dagger\bigl(d_1^\dagger B\bigr) \;=\; \phaseRingCount\,(d_0^\dagger d_1^\dagger n) \;\equiv\; 0.
  \end{equation}
\end{enumerate}
\end{theorem}

\begin{proof}
Равенство~\eqref{eq:discrete-plaquette-flux} следует из того, что сумма разностей $\varphi(y)-\varphi(x)$
вдоль замкнутого контура $\partial f$ телескопически сокращается: $\bigl(d_1(d_0\varphi)\bigr)_f = 0$,
поэтому $(d_1 A)_f = -\phaseRingCount(d_1 k)_f = \phaseRingCount n_f$.
Равенство~\eqref{eq:discrete-faraday} есть $d_1(d_0\,\Delta_t\varphi)=0$.
В~\eqref{eq:discrete-gauss-b} каждое из $24$ рёбер кубооктаэдра принадлежит ровно двум граням из $8\triangle\cup 6\square$
с противоположными ориентациями обхода, поэтому $d_2\circ d_1\equiv 0$.
Подстановка~\eqref{eq:discrete-edge-connection} и~\eqref{eq:discrete-plaquette-flux} в $d_0^\dagger E$ и $d_1^\dagger B-\Delta_t E$
даёт~\eqref{eq:discrete-ampere-gauss}, а применение $d_0^\dagger$ к $d_1^\dagger B$ даёт~\eqref{eq:discrete-charge-conservation}
в силу $(d_1 d_0)^\dagger = 0$.
\end{proof}

\begin{corollary}[Единство природы и количества заряда: частное и остаток одного деления на $\phaseRingCount = 512$]
\label{cor:charge-origin-and-count-quotient-remainder}
В разложении реберной разности фаз~\eqref{eq:discrete-edge-connection}
\[
  (d_0\varphi)_{xy} \;=\; \varphi(y) - \varphi(x) \;=\; \phaseRingCount\,k_{xy} + A_{xy},
  \qquad
  A_{xy}\in(-\Phi_{\mathrm{Pauli}},\,+\Phi_{\mathrm{Pauli}}],\quad k_{xy}\in\{-1,0,+1\}
\]
скачок ветви $k_{xy}$ и связность $A_{xy}$ не являются двумя независимыми полями: по теореме об евклидовом делении
с остатком на модуль кольца $\phaseRingCount = 2^{1+|\triangle|} = 512$ со сдвигом в главный интервал $(-256,256]$
они суть \textbf{неполное частное} и \textbf{остаток} одной и той же целочисленной разности $(d_0\varphi)_{xy}$:
\begin{equation}
\label{eq:quotient-remainder-512}
  k_{xy} \;=\; \left\lfloor \frac{(d_0\varphi)_{xy} + \Phi_{\mathrm{Pauli}} - 1}{\phaseRingCount} \right\rfloor,
  \qquad
  A_{xy} \;=\; (d_0\varphi)_{xy} - \phaseRingCount\,k_{xy}.
\end{equation}
Отсюда без какого-либо отдельного постулата однозначно следуют не только природа заряда («откуда заряд»),
но и его количественная доля по отношению к излучению («сколько его»):
\begin{enumerate}[label=(\roman*)]
\item \textbf{Частное есть материя (топологический заряд), остаток есть излучение (фотонная ванна):}
  пока на ребре разность фаз остаётся внутри главного интервала ($|(d_0\varphi)_{xy}|\le\Phi_{\mathrm{Pauli}}=256$),
  неполное частное равно нулю ($k_{xy}=0 \implies n_f = -(d_1 k)_f = 0$), и вся надпороговая активность ворот
  $|\Phi|\ge\Delta\varphi_{\mathrm{disc}}=41$ остаётся в остатке $A_{xy}$ в виде безысточниковых квантов излучения
  $n_E = \lfloor|\Phi|/41\rfloor$. Как только под действием градиента ворот $(d_0\Phi)_{xy}$ реберная разность
  перешагивает полуоборот Паули $\Phi_{\mathrm{Pauli}}=256$, деление~\eqref{eq:quotient-remainder-512} перебрасывает
  единицу из остатка $A_{xy}$ в частное ($\Delta_t k_{xy}=\pm 1$), рождая на примыкающих гранях целочисленный
  плакетный заряд $n_f = -(d_1 k)_f = (d_1 A)_f/\phaseRingCount \in \mathbb{Z}\setminus\{0\}$.
\item \textbf{Рождение ровно $2N_\varphi = 26$ вихревых плакетов из безвихревого начального вектора $\mathrm{IV}(3)=263$ на такте $t=2$:}
  в начальном состоянии $\symCapitalPsi^0 = \mathrm{IV}(3) = 2^{2^3} + (2^3-1) = 263$ (\cref{thm:catalan-uniqueness-d3})
  фаза задана узловой $0$-коцепью $\varphi^0(x)$, для которой $d_1(d_0\varphi^0)\equiv 0$, поэтому при $t=0$ и $t=1$
  вихревых плакетов нет во всём пространстве ($n_f(t=0)\equiv 0$, $n_f(t=1)\equiv 0$ --- заряды не всаживаются
  в начальное условие). На такте $t=2$ (двухшаговое замыкание контура плакета на первой сфере кубооктаэдра~$S_1$)
  переброс частного $\Delta_t k_{xy}=\pm 1$ через порог Паули $256$ тиков на границе трёхгранного угла~$K_3$
  рождает из этого безвихревого состояния ровно
  \begin{equation}
  \label{eq:t2-26-vortex-plaquettes}
    \sum_{f}\bigl|n_f(t=2)\bigr| \;=\; 2\,N_\varphi \;=\; 2\bigl(1 + |N_{12}|\bigr) \;=\; 26
  \end{equation}
  вихревых плакетов --- полный удвоенный по спину кластер кубооктаэдра.
\item \textbf{Однозначность барион-фотонного отношения $\eta_B$ как отношения выходов частного и остатка деления~\eqref{eq:quotient-remainder-512}:}
  поскольку частное $k_{xy}$ и остаток $A_{xy}$ порождаются одним и тем же делением $(d_0\varphi)_{xy}$ на $\phaseRingCount=512$,
  отношение числа выживших барионных узлов ($n_f\ne 0$) к числу фотонных квантов ($n_E\ge 1,\ n_f=0$) есть доля
  переполнений главного интервала в полном потоке кубооктаэдра: на $2N_\varphi = 26$ вихревых плакетов~\eqref{eq:t2-26-vortex-plaquettes}
  CP-нечётная мнимая компонента гауссовой факторизации кулоновского канала $\hat M = d_{\mathrm{in}}(\square)^2 + \mathrm{mod\_bits}^2 = 4^2+9^2 = 97$
  составляет $\bar\eta = \mathrm{mod\_bits}/(2N_\varphi) = 9/26$, доля внешних лучей рёберного канала $E=(2,5,5)$
  на $|\triangle|=2^d=8$ октантных ступеней равна $d_{\mathrm{out}}(E)/|\triangle| = 5/8$, а прохождение через
  $d_{\mathrm{in}}(\square)=4$ квадранта квадратной грани даёт фактор $\alpha_{\mathrm{fs}}^{d_{\mathrm{in}}(\square)}=\alpha_{\mathrm{fs}}^4$:
  \begin{equation}
  \label{eq:eta-b-from-quotient-remainder}
    \eta_B \;=\; \frac{\mathrm{mod\_bits}}{2\,N_\varphi}\cdot\frac{d_{\mathrm{out}}(E)}{|\triangle|}\,\alpha_{\mathrm{fs}}^{d_{\mathrm{in}}(\square)}
    \;=\; \frac{9}{26}\cdot\frac{5}{8}\,\alpha_{\mathrm{fs}}^4
    \;=\; \frac{45}{208}\,\alpha_{\mathrm{fs}}^4 \;\approx\; 6{,}135\times 10^{-10}.
  \end{equation}
\end{enumerate}
\end{corollary}

\begin{remark}[Единая геометрия всех взаимодействий в одних воротах на кубооктаэдре $N_{12}$]
\label{rem:unified-interactions-single-gate}
Теоремы~\ref{thm:holonomy-split-continuity} и~\ref{thm:discrete-maxwell-sourceless} показывают, что на~$\Mlayer$
нет ни отдельных «частиц-источников», ни отдельных «полей-переносчиков»: все взаимодействия являются
разными проекциями \textbf{одних и тех же ворот} $\Phi$ на одном девятибитном регистре $(\mu,\varphi)$
и одном клеточном комплексе кубооктаэдра ($12$ вершин, $24$ ребра, $8\triangle + 6\square$ граней):
\begin{enumerate}[label=(\arabic*)]
\item \textbf{Квантовая механика (Шрёдингер, Гейзенберг, Паули):}
  чётный дефицит голономии $\delta\zeta_{\mathrm{real}}$~\eqref{eq:even-holonomy-phase} даёт кинетический
  член фазы и дисперсию волн; нечётная часть $\zeta_{\mathrm{imag}}$~\eqref{eq:odd-holonomy-continuity}
  даёт уравнение непрерывности без источника; порог $\Delta\varphi_{\mathrm{disc}}=41$ тик задаёт мёртвую зону
  Гейзенберга, а полуоборот $\Phi_{\mathrm{Pauli}}=256$ тиков при насыщении потолка $\mu_{\max}=64$ --- запрет Паули.
\item \textbf{Электромагнетизм ($U(1)$) и постоянная тонкой структуры $\alpha_{\mathrm{fs}}$:}
  абелева связность $A_{xy}$~\eqref{eq:discrete-edge-connection} на $12$ связях и поток $B_f=\phaseRingCount n_f$~\eqref{eq:discrete-plaquette-flux}
  через грани кубооктаэдра подчиняются точным дискретным уравнениям Максвелла~\eqref{eq:discrete-faraday}--\eqref{eq:discrete-ampere-gauss},
  а константа связи $\alpha_{\mathrm{fs}}^{-1}=\pi^2(|N_{12}|+\kappa)+\dots\approx 137{,}036$ (гл.~\ref{ch:alpha})
  есть доля одного реберного канала в полном фазовом бюджете кубооктаэдра при $\hat M = 1+|N_{12}|\,|\triangle| = 97$.
\item \textbf{Слабое и сильное взаимодействия ($SU(2)\times SU(3)$):}
  двухкомпонентный спинор $(z_1,z_2)$ несёт точную локальную группу $SU(2)$ на звезде, а распределение
  неабелева потока по оболочке кубооктаэдра распадается по двум типам её двумерных граней:
  $|\triangle|=8$ треугольных граней (восемь компонент глюонного октета, $\dim\mathfrak{su}(3)=8$)
  и $|\square|=6$ квадратных граней (три координатные плоскости $\times$ две ориентации).
\item \textbf{Гравитация (локальное замедление такта, знак градиента фазовой скорости и преломление Редже):}
  вес ячейки $q(x)=\mu^2(x)=\rho_{\mathrm{vac}}(1+\varepsilon_x)$ стоит в знаменателе тех же самых ворот
  $D(x) := \zeta_{\mathrm{real}}(x)+\mu^2(x)+K_P$ (\cref{ch:evolution}).
  Прямой замер на носителе ($24^3$, канонический 3D-вихрь \texttt{VORTEX\_P} на кипящем океане)
  даёт монотонный рост знаменателя~$D(r)$ от периферии к центру одетого дефекта:
  в далёком вакууме $q_{\mathrm{vac}}=2$, $\zeta_{\mathrm{real}}=16$, $D_{\mathrm{vac}}=82\ \text{кванта}^2$;
  на сфере~$S_3$ ($k=3$) $\langle q\rangle=158$, $\langle D\rangle=2077$;
  на~$S_2$ ($k=2$) $\langle q\rangle=826{,}4$, $\langle D\rangle=9426{,}7$;
  на одевающей сфере~$S_1$ ($k=1$, $R_{\mathrm{dress}}=1\cdot dl$) $\langle q\rangle=2499{,}3$, $\langle D\rangle=24138{,}0$;
  в ядре ($k=0$) $q=3929$, $D=32493\ \text{кванта}^2$.
  Поскольку приращение фазы пробной волны за один такт~$dt$ обратно пропорционально знаменателю ворот,
  \[
    \Delta\Phi_{\mathrm{test}}(r) \;\propto\; \frac{1}{\zeta_{\mathrm{real}}(r)+\mu^2(r)+K_P} \;=\; \frac{1}{D_{\mathrm{vac}}\bigl(1+\varepsilon_{\mathrm{eff}}(r)\bigr)},
    \qquad
    \frac{d D(r)}{d r} < 0,
  \]
  угловая скорость набега фазы при приближении к массивному узлу ($r\to 0$) строго \textbf{убывает}.
  Поворот волнового фронта при этом происходит \textbf{строго локально на одном ребре за один такт~$dt$},
  без обращения к глобальному вариационному принципу Ферма. Действительно, для любых двух соседних узлов
  поперечного ребра $e_\perp = (x_{\mathrm{out}}\to x_{\mathrm{in}})$ (где $x_{\mathrm{in}}$ лежит на один шаг~$dh$
  ближе к массивному центру, так что $D(x_{\mathrm{in}}) > D(x_{\mathrm{out}})$) приращение реберной разности фаз
  $A_\perp = \mathrm{wrap}(\varphi(x_{\mathrm{in}})-\varphi(x_{\mathrm{out}}))$ за один такт~$dt$ равно когранице
  локального шага ворот~$\Phi$ (частный случай дискретного закона Фарадея~\eqref{eq:discrete-faraday} при $B_f=0$):
  \begin{equation}
  \label{eq:local-edge-refraction}
    \Delta_t A_\perp(x_{\mathrm{out}}\to x_{\mathrm{in}})
    \;=\;
    (d_0\Phi)_{x_{\mathrm{out}}\to x_{\mathrm{in}}}
    \;=\;
    \Phi(x_{\mathrm{in}}) - \Phi(x_{\mathrm{out}})
    \;<\; 0.
  \end{equation}
  Точка фронта~$x_{\mathrm{in}}$, ближайшая к массивному узлу, за один такт поворачивает фазу на меньшее число тиков,
  чем соседняя внешняя точка~$x_{\mathrm{out}}$ того же ребра, поэтому дискретный волновой вектор фронта на каждом
  такте получает поперечное приращение вдоль локального градиента знаменателя ворот $\nabla_{\mathrm{disc}} D(x)$,
  заворачивая фронт \textbf{к массивному центру}, а не от него.
  Континуальный показатель преломления вакуума $n(r) := D(r)/D_{\mathrm{vac}} = 1 + \varepsilon_{\mathrm{eff}}(r) > 1$
  ($dn/dr < 0$), принцип Ферма ($\delta\int n\,ds=0$) и угол гравитационного отклонения $\theta_{\mathrm{def}}=4GM/(c^2 b)$
  (\cref{rem:channel-gr-open-in-book}) являются лишь длинноволновой суммой этих одношаговых реберных поворотов~\eqref{eq:local-edge-refraction} на~$\Tlayer$.
\end{enumerate}
Таким образом, гравитация (знаменатель ворот $\mu^2$), квантовая дисперсия (чётная часть числителя $\delta\zeta_{\mathrm{real}}$),
поток плотности (нечётная часть числителя $\zeta_{\mathrm{imag}}=-\operatorname{div}_{\mathrm{disc}}\mathbf{j}$),
калибровочные поля (циркуляция $d_1 A$ по граням $8\triangle\cup 6\square$) и материя (узлы с топологическим
числом $n_f\ne 0$) не постулируются порознь, а суть пять граней одного арифметического тождества на~$\Mlayer$.
\end{remark}

\begin{theorem}[Закрытая алгебра реберных проекций на гранях кубооктаэдра без тригонометрических функций]
\label{thm:trig-free-face-projections}
На~$\Mlayer$ носитель не использует ни трансцендентных функций $\cos,\sin$, ни числа~$\pi$, ни таблиц
с плавающей запятой: чётная и нечётная проекции реберной разности фаз $A_{xy} = \mathrm{wrap}(\varphi(y)-\varphi(x))$
полностью определяются комбинаторными инвариантами граней одношаговой оболочки кубооктаэдра в размерности~$d=3$:
\begin{equation}
\label{eq:cubocta-face-invariants}
  |N_{12}| = 2d(d-1) = 12,\quad
  |\triangle| = 2^d = 8,\quad
  |\square| = 2d = 6,\quad
  d_{\mathrm{in}}(\square) = \frac{|N_{12}|}{d} = 2(d-1) = 4.
\end{equation}
Пусть $\mathrm{mod\_bits} = 1 + |\triangle| = d^2 = 9$, $\mathrm{frac\_bits} = \mathrm{mod\_bits} - d = |\square| = 6$,
$\phaseRingCount = 2^{\mathrm{mod\_bits}} = 2^{1+|\triangle|} = 512$ и $K_P = 2^{\mathrm{frac\_bits}} = 2^{|\square|} = 64$.
Экваториальная квадратная грань~$\square$ с $d_{\mathrm{in}}(\square)=4$ вершинами делит фазовое кольцо~$\mathbb{Z}_{\phaseRingCount}$
на $d_{\mathrm{in}}(\square)$ равных квадрантов длины
\begin{equation}
\label{eq:quadrant-pauli-closed}
\begin{aligned}
  Q_\square &\;:=\; \frac{\phaseRingCount}{d_{\mathrm{in}}(\square)} \;=\; 2^{|\triangle|-1} \;=\; 2\,K_P \;=\; 128,\\
  \Phi_{\mathrm{Pauli}} &\;=\; \frac{\phaseRingCount}{2} \;=\; 2\,Q_\square \;=\; 2^{|\triangle|} \;=\; d_{\mathrm{in}}(\square)\,K_P \;=\; 256.
\end{aligned}
\end{equation}
Тогда для каждого направленного ребра $e=(x\to y)$ со свёрнутой разностью фаз $A_{xy}\in[-\Phi_{\mathrm{Pauli}},\,+\Phi_{\mathrm{Pauli}}]$:
\begin{enumerate}[label=(\roman*)]
\item \textbf{Чётная проекция ребра (когерентность) и чётное натяжение ребра:}
  \begin{equation}
  \label{eq:discrete-cos-closed}
  \begin{aligned}
    C(A_{xy}) &\;:=\; Q_\square - |A_{xy}| \;=\; \frac{\phaseRingCount}{d_{\mathrm{in}}(\square)} - |A_{xy}| \;\in\; [-Q_\square,\,+Q_\square],\\
    T(A_{xy}) &\;:=\; Q_\square - C(A_{xy}) \;=\; |A_{xy}| \;\ge\; 0.
  \end{aligned}
  \end{equation}
\item \textbf{Нечётная проекция ребра (поток, сдвинутый на квадрант $Q_\square$):}
  \begin{equation}
  \label{eq:discrete-sin-closed}
  \begin{aligned}
    S(A_{xy}) &\;:=\; C\bigl(\mathrm{wrap}(A_{xy} - Q_\square)\bigr)
    \;=\; \operatorname{sgn}(A_{xy})\cdot\Bigl(Q_\square - \bigl||A_{xy}| - Q_\square\bigr|\Bigr)\\
    &\;=\; -S(A_{yx}) \;\in\; [-Q_\square,\,+Q_\square].
  \end{aligned}
  \end{equation}
\item \textbf{Метрическое тождество на квадратной грани ($L^1$-инвариант) и точный пифагоров ротатор ($L^2$-инвариант без $\pi$):}
  для всех $A_{xy}\in\mathbb{Z}_{\phaseRingCount}$ сумма модулей чётной и нечётной проекций тождественно постоянна:
  \begin{equation}
  \label{eq:l1-face-identity}
    \bigl|C(A_{xy})\bigr| + \bigl|S(A_{xy})\bigr| \;\equiv\; Q_\square \;=\; \frac{\phaseRingCount}{d_{\mathrm{in}}(\square)} \;=\; 2\,K_P,
  \end{equation}
  а целочисленная пифагорова тройка Кэли на квадранте квадратной грани
  \begin{equation}
  \label{eq:pythagorean-rotator-int}
    P_x(A) := \operatorname{sgn}(C(A))\bigl(Q_\square^2 - S(A)^2\bigr),\quad
    P_y(A) := 2\,Q_\square\,S(A),\quad
    P_r(A) := Q_\square^2 + S(A)^2
  \end{equation}
  удовлетворяет точному основному тригонометрическому тождеству в целых числах:
  \begin{equation}
  \label{eq:pythagorean-identity-int}
    P_x(A)^2 + P_y(A)^2 \;\equiv\; P_r(A)^2
    \quad\Longleftrightarrow\quad
    \left(\frac{Q_\square^2 - S^2}{Q_\square^2 + S^2}\right)^{\!2} + \left(\frac{2\,Q_\square\,S}{Q_\square^2 + S^2}\right)^{\!2} \;\equiv\; 1
    \quad\forall\,A\in\mathbb{Z}_{\phaseRingCount}.
  \end{equation}
  В частности, на пороге Гейзенберга $A = \Delta\varphi_{\mathrm{disc}} = 41$ при $Q_\square = 128$ ротатор~\eqref{eq:pythagorean-rotator-int}
  даёт точную пифагорову тройку Гейзенберга $(P_x, P_y, P_r) = (14703,\,10496,\,18065)$, где $14703^2 + 10496^2 = 18065^2$.
\item \textbf{Выход к гармоническим $\cos$ и $\sin$ на~$\Tlayer$:}
  при мягком чтении~$\mathcal{B}$ с отношением граней $\caKappaGeom^2 = |\square|/|N_{12}| = 1/(d-1) = 1/2$ ($\caKappaGeom = 1/\sqrt{2}$)
  и $|\triangle|/|\square| = 2^{d-1}/d = 4/3$ многократная свёртка кусочно-линейных реберных проекций
  $(C,S)/Q_\square$ по шагам оболочки кубооктаэдра переводит $L^1$-ромб~\eqref{eq:l1-face-identity}
  в гладкую $L^2$-окружность $\cos^2\varphi + \sin^2\varphi = 1$ на~$\Tlayer$ (\cref{thm:t-cl}).
\end{enumerate}
\end{theorem}

\begin{proof}
Из $|N_{12}|=2d(d-1)$, $|\triangle|=2^d$ и $|\square|=2d$ при $d=3$ следуют~\eqref{eq:cubocta-face-invariants}
и~\eqref{eq:quadrant-pauli-closed}, так как $d_{\mathrm{in}}(\square)=4=2^2$ и $\phaseRingCount/d_{\mathrm{in}}(\square)=2^{9-2}=2^{8-1}=2^{|\triangle|-1}=2K_P$.
При $|A_{xy}|\in[0,Q_\square]$ имеем $C(A_{xy}) = Q_\square - |A_{xy}|\ge 0$ и $|S(A_{xy})| = Q_\square - (Q_\square - |A_{xy}|) = |A_{xy}|$,
откуда $|C(A_{xy})|+|S(A_{xy})| = (Q_\square - |A_{xy}|) + |A_{xy}| = Q_\square$.
При $|A_{xy}|\in(Q_\square, 2Q_\square]$ имеем $C(A_{xy}) = Q_\square - |A_{xy}| < 0$, так что $|C(A_{xy})| = |A_{xy}| - Q_\square$,
а $|S(A_{xy})| = Q_\square - (|A_{xy}| - Q_\square) = 2Q_\square - |A_{xy}|$, и их сумма снова равна
$(|A_{xy}| - Q_\square) + (2Q_\square - |A_{xy}|) = Q_\square$, что доказывает~\eqref{eq:l1-face-identity}.
Антисимметрия $S(A_{yx}) = S(-A_{xy}) = -S(A_{xy})$ влечёт $\sum_{x\in\carrierLattice}\sum_{y\in N_{12}(x)} S(A_{xy}) \equiv 0$.
\end{proof}

\begin{theorem}[Диофантова единственность $d=3$ ($1+2^d=d^2$), единственность начального вектора $\mathrm{IV}(3)=263$ и полная замкнутость траектории $\symCapitalPsi^t = g^t(\mathrm{IV}(3))$]
\label{thm:catalan-uniqueness-d3}
Пусть локальный обратимый клеточный автомат в пространственной размерности $d\in\mathbb{N}$ ($d\ge 1$)
имеет оболочку ближайших соседей с $|N(d)| = 2d(d-1)$ вершинами, $|\triangle(d)| = 2^d$ октантными гранями,
$|\square(d)| = 2d$ координатными гранями и $d_{\mathrm{in}}(\square) = |N(d)|/d = 2(d-1)$.
Потребуем выполнения трёх арифметических условий замкнутости носителя:
\begin{enumerate}[label=(\roman*)]
\item \textbf{Замок битового регистра:} число «центральный узел плюс октанты» совпадает с квадратом размерности:
  \begin{equation}
  \label{eq:catalan-lock}
    \mathrm{mod\_bits}(d) \;:=\; 1 + 2^d \;=\; d^2.
  \end{equation}
\item \textbf{Замок целочисленности сектора Гейзенберга на кольце:} при объёме кольца $\phaseRingCount(d) = 2^{1+2^d}$
  и размере одношагового кластера $N_\varphi(d) = 1 + 2d(d-1)$ отношение
  \begin{equation}
  \label{eq:heisenberg-divisibility-lock}
    \Delta\varphi_{\mathrm{disc}}(d) \;:=\; \frac{2^{1+2^d} + \bigl(1 + 2d(d-1)\bigr) + 2^d}{1 + 2d(d-1)} \;\in\; \mathbb{Z}
  \end{equation}
  делится нацело без остатка.
\item \textbf{Замок гауссовой факторизации кулоновского канала в $\mathbb{Z}[i]$:} площадь «звезда $\times$ октанты плюс ядро»
  раскладывается в сумму квадратов стороны квадратной грани и разрядности регистра:
  \begin{equation}
  \label{eq:gaussian-coulomb-lock}
    \hat M(d) \;:=\; 1 + 2d(d-1)\cdot 2^d \;=\; \bigl(2(d-1)\bigr)^2 + \bigl(1 + 2^d\bigr)^2.
  \end{equation}
\end{enumerate}
Тогда каждое из диофантовых уравнений~\eqref{eq:catalan-lock} и~\eqref{eq:gaussian-coulomb-lock} в отдельности, как и вся система~\eqref{eq:catalan-lock}--\eqref{eq:gaussian-coulomb-lock}, во всём множестве натуральных чисел $d\in\mathbb{N}$
имеет \textbf{одно и только одно решение $d = 3$}, из которого без единого свободного параметра однозначно следуют как все инварианты носителя и закона~$g$:
\[
  |N_{12}|=12,\;\; |\triangle|=8,\;\; |\square|=6,\;\; d_{\mathrm{in}}(\square)=4,\;\; \caKappaGeom^2=\frac{1}{2}\ \Bigl(\caKappaGeom=\frac{1}{\sqrt{2}}\Bigr),\;\;
  \mathrm{mod\_bits}=9,\;\; \mathrm{frac\_bits}=6,
\]
\[
  \phaseRingCount=512,\;\; K_P=64,\;\; Q_\square=128,\;\; \Phi_{\mathrm{Pauli}}=256,\;\;
  N_\varphi=13,\;\; \Delta\varphi_{\mathrm{disc}}=41,\;\; \hat M = 4^2 + 9^2 = 97,
\]
так и \textbf{единственный $9$-битный начальный вектор нуклеации Вселенной} $\symCapitalPsi^{0} = \mathrm{IV}(d)$ на фоне кипящего океана~$A_5$:
\begin{equation}
\label{eq:unique-initial-vector-d3}
  \mathrm{IV}(d) \;:=\; 2^{2^d} + \bigl(2^d - 1\bigr) \;=\; 2^{8} + \bigl(2^{3} - 1\bigr) \;=\; 256 + 7 \;=\; \mathbf{263} \;=\; \texttt{0b1\_00000111},
\end{equation}
где старший бит $2^{|\triangle|}=256$ есть единичный бит хиральности центрального ядра ($b_{\mathrm{chir}}=1$),
а младшие $|\triangle|=8$ бит несут маску трёхгранного барионного угла~$K_3$ по $d=3$ координатным измерениям ($2^d-1=7$),
с радиусом ядра $R_{\mathrm{IV}}(d) = |\square| = 2d = 6\,dh$, амплитудой ядра $\mu_{\mathrm{IV}}(d) = N_\varphi = 1 + 2d(d-1) = 13$ квантов
и фазовым сдвигом октантов $\Delta\varphi_{\mathrm{IV}}(d) = d\cdot\Delta\varphi_{\mathrm{disc}}(d) = 3\times 41 = 123$ тика кольца~$\mathbb{Z}_{512}$.

Следовательно, ни геометрия носителя, ни закон эволюции~$g$, ни фундаментальные константы, ни начальное состояние Вселенной~$\symCapitalPsi^{0}$
не содержат свободных параметров: вся физическая история есть единственная детерминированная целочисленная траектория
\begin{equation}
\label{eq:closed-universe-trajectory}
  \symCapitalPsi^{t} \;=\; g^{t}\bigl(\mathrm{IV}(3)\bigr),
  \qquad
  d = 3,\quad \mathrm{IV}(3) = 263,\quad t\in\mathbb{Z}.
\end{equation}
\end{theorem}

\begin{proof}
\emph{(i)} Проверим $d\in\{1,2,3\}$ прямой подстановкой в~\eqref{eq:catalan-lock}:
при $d=1$ имеем $1+2^1 = 3 \ne 1^2 = 1$; при $d=2$ имеем $1+2^2 = 5 \ne 2^2 = 4$; при $d=3$ выполняется $1+2^3 = 9 = 3^2$.
Докажем по индукции, что для всех целых $d\ge 4$ выполняется строгое неравенство $2^d + 1 > d^2$.
База индукции ($d=4$): $2^4 + 1 = 17 > 16 = 4^2$.
Индуктивный переход ($d\to d+1$ при $d\ge 4$): пусть $2^d + 1 > d^2$. Умножая обе части на $2$ и вычитая $1$, получаем
\[
  2^{d+1} + 1 \;=\; 2(2^d + 1) - 1 \;>\; 2d^2 - 1.
\]
Вычислим разность между $2d^2 - 1$ и $(d+1)^2 = d^2 + 2d + 1$:
\[
  (2d^2 - 1) - (d^2 + 2d + 1) \;=\; d^2 - 2d - 2 \;=\; (d-1)^2 - 3.
\]
Поскольку $d\ge 4$, имеем $d-1 \ge 3$, откуда $(d-1)^2 - 3 \ge 3^2 - 3 = 6 > 0$, то есть $2d^2 - 1 > (d+1)^2$,
и по транзитивности $2^{d+1} + 1 > (d+1)^2$. Следовательно, при $d\ge 4$ уравнение~\eqref{eq:catalan-lock} решений не имеет,
и $d=3$ есть единственное решение~\eqref{eq:catalan-lock} во всём~$\mathbb{N}$.

\emph{(iii)} Раскроем квадрат $(1+2^d)^2 = 1 + 2\cdot 2^d + 2^{2d}$ в правой части~\eqref{eq:gaussian-coulomb-lock} и вычтем $1$ из обеих частей:
\[
  2d(d-1)\cdot 2^d \;=\; 4(d-1)^2 + 2\cdot 2^d + 2^{2d}
  \quad\Longleftrightarrow\quad
  2^d\bigl(2d(d-1) - 2 - 2^d\bigr) \;=\; 4(d-1)^2.
\]
Проверим $d\in\{1,2,3,4,5\}$ прямой подстановкой:
при $d=1$ левая часть равна $2^1(0-2-2) = -8 \ne 0 = 4(1-1)^2$;
при $d=2$ левая часть равна $2^2(4-2-4) = -8 \ne 4 = 4(2-1)^2$;
при $d=3$ левая часть равна $2^3(12-2-8) = 8\times 2 = 16 = 4(3-1)^2$ (равенство выполняется: $1+12\cdot 8 = 97 = 4^2+9^2$);
при $d=4$ левая часть равна $2^4(24-2-16) = 16\times 6 = 96 \ne 36 = 4(4-1)^2$;
при $d=5$ левая часть равна $2^5(40-2-32) = 32\times 6 = 192 \ne 64 = 4(5-1)^2$.
Докажем по индукции, что для всех целых $d\ge 6$ выполняется строгое неравенство $2^d > 2d(d-1) - 2$.
База индукции ($d=6$): $2^6 = 64 > 58 = 2\cdot 6\cdot 5 - 2$.
Индуктивный переход ($d\to d+1$ при $d\ge 6$): пусть $2^d > 2d(d-1) - 2$. Умножая на $2$, получаем $2^{d+1} > 4d(d-1) - 4$.
Разность между $4d(d-1) - 4$ и $2(d+1)d - 2 = 2d^2 + 2d - 2$ равна
\[
  (4d^2 - 4d - 4) - (2d^2 + 2d - 2) \;=\; 2d^2 - 6d - 2 \;=\; 2d(d-3) - 2 \;\ge\; 12\cdot 3 - 2 \;=\; 34 \;>\; 0,
\]
откуда $2^{d+1} > 2(d+1)d - 2$. Таким образом, для всех $d\ge 6$ левая часть $2^d(2d(d-1)-2-2^d) < 0$, тогда как правая часть $4(d-1)^2 > 0$.
Следовательно, $d=3$ есть единственное решение~\eqref{eq:gaussian-coulomb-lock} во всём~$\mathbb{N}$.

\emph{(ii)} Запишем отношение~\eqref{eq:heisenberg-divisibility-lock} в виде
\[
  \Delta\varphi_{\mathrm{disc}}(d) \;=\; 1 + \frac{2^d\bigl(2^{2^d - d + 1} + 1\bigr)}{1 + 2d(d-1)}.
\]
При единственном корне $d=3$ уравнений~\eqref{eq:catalan-lock} и~\eqref{eq:gaussian-coulomb-lock} имеем $2^d - d + 1 = 8 - 3 + 1 = 6$,
$2^6 + 1 = 65 = 5\times 13$ и $N_\varphi(3) = 1 + 2\cdot 3\cdot 2 = 13$, откуда
\[
  \Delta\varphi_{\mathrm{disc}}(3) \;=\; 1 + \frac{8\times 65}{13} \;=\; 1 + 40 \;=\; 41 \;\in\; \mathbb{Z}.
\]
Таким образом, $d=3$ --- единственное решение системы~\eqref{eq:catalan-lock}--\eqref{eq:gaussian-coulomb-lock} в~$\mathbb{N}$.
\end{proof}

\begin{remark}[Откуда берётся $2\pi$ в $\planckHbar = h/(2\pi)$ и что такое само число $\pi$ на решётке]
\label{rem:hbar-is-t-artifact}
На~$\Mlayer$ элементарным атомом действия является один тик фазового кольца $s_0 = 1$, полный оборот кольца
состоит ровно из $\phaseRingCount = 2^{1+|\triangle|} = 512$ тиков ($h = \phaseRingCount\,s_0 = 512\,s_0$),
а порог Гейзенберга равен целому числу $\Delta\varphi_{\mathrm{disc}} = (512+13+8)/13 = 41$ тик ($41\,s_0$).
Ни числа~$2\pi$, ни «приведённой постоянной Планка» $\planckHbar$ на~$\Mlayer$ не существует:
$\planckHbar := h/(2\pi) = (512/(2\pi))\,s_0 \approx 81{,}487\,s_0$ (и $\planckHbar/2 \approx 40{,}744\,s_0$)
возникает только на~$\Tlayer$, когда наблюдатель переводит счёт целых оборотов $512$-угольника в радианную меру~$2\pi$.

Более того, само число~$\pi$ не является в основании геометрии «отношением длины идеальной окружности к диаметру»
(на~$\Mlayer$ нет континуальной окружности, а на одном шаге $N=1$ отношение полуоборота кольца $\Phi_{\mathrm{Pauli}}=256$
к двум секторам Гейзенберга $2\Delta\varphi_{\mathrm{disc}}=82$ есть рациональная дробь $\pi_1 = Q_\square/\Delta\varphi_{\mathrm{disc}} = 128/41 \approx 3{,}122$).
Трансцендентное число~$\pi$ рождается при подъёме $\Mlayer\to\Tlayer$ как \textbf{асимптотическая комбинаторная плотность
возврата замкнутых путей} на дискретной решётке при числе шагов $N\to\infty$ (формула Валлиса--Стирлинга для числа
возвращающихся петель $\binom{2N}{N} \sim 2^{2N}/\sqrt{\pi N}$ из $2^{2N}$ битовых маршрутов, \cref{thm:t-cl}):
\[
  \pi \;=\; \lim_{N\to\infty} \frac{1}{N}\left(\frac{2^{2N}}{\dbinom{2N}{N}}\right)^{\!2}
  \;=\; 2\prod_{k=1}^{\infty}\frac{4k^2}{4k^2-1}.
\]
Каждый следующий знак числа~$\pi$ кодирует следующий порядок комбинаторных корреляций многошаговых замкнутых
путей на кубооктаэдральной решётке $g^t(d=3)$.
Точно так же элементарный такт времени~$dt$ есть абсолютная единица времени носителя, а длительность макроскопической
секунды СИ есть фиксированное число тактов $\mathcal{N}_{\mathrm{sec}}(d=3)\,dt$, вычисляемое через
комбинаторную лестницу масс и частоту перехода водорода без участия гравитационной постоянной~$G$.
\end{remark}

\section{Вся физика в целых числах \texorpdfstring{$\mathbb{Z}$}{Z} и происхождение континуальных теорий на~\texorpdfstring{$\Tlayer$}{T}: Шрёдингер, Навье--Стокс, Янг--Миллс со щелью массы и Эйнштейн}
\label{sec:z-to-continuum-theories}

Из теорем~\ref{thm:holonomy-split-continuity}--\ref{thm:catalan-uniqueness-d3} следует, что на фундаментальном носителе~$\Mlayer$
природа не использует ни вещественных чисел~$\mathbb{R}$, ни дифференциального исчисления, ни бесконечномерных функциональных интегралов.
Вся точная микрофизика исчерпывается четырьмя арифметическими правилами в целых числах~$\mathbb{Z}$, доступными прямому счёту:
\begin{enumerate}[label=\textbf{Z\arabic*.},leftmargin=*]
\item \textbf{Трёхмерность и геометрия упаковки из одного уравнения $1+2^d=d^2$:}
  единственный корень $d=3$ фиксирует кубооктаэдр с $|N_{12}|=2d(d-1)=12$ соседями, $|\triangle|=2^d=8$ треугольными
  и $|\square|=2d=6$ квадратными гранями ($d_{\mathrm{in}}(\square)=4$), разрядность слова $\mathrm{mod\_bits}=1+2^d=9$
  ($\phaseRingCount=2^9=512$) и мантиссу амплитуды $\mathrm{frac\_bits}=9-3=6$ ($\mu\in\{1,\dots,64\}$).
\item \textbf{Часы на $512$ делений, порог Гейзенберга $41$ и зеркало Паули $256$:}
  на каждом ребре $(x\to y)$ работают кусочно-линейные проекции квадратной грани $C(A)=128-|A|$ и $S(A)=C(\mathrm{wrap}(A-128))$
  с тождеством $|C(A)|+|S(A)|\equiv 128$; ворота $\Phi\in\mathbb{Z}_{512}$ открываются только при $|\Phi|\ge\Delta\varphi_{\mathrm{disc}}=(512+13+8)/13=41$ тик,
  а при параллельном перекрытии двух спиноров в ячейке срабатывает полуоборот $\Phi_{\mathrm{Pauli}}=256$ тиков.
\item \textbf{Свет и материя как остаток и частное одного деления на $512$:}
  при евклидовом делении реберной разности фаз $(d_0\varphi)_{xy} = 512\,k_{xy} + A_{xy}$ со сдвигом в $(-256,256]$
  остаток $A_{xy}$ переносит безысточниковые волны излучения ($n_E = \lfloor|\Phi|/41\rfloor$), а кограница неполного частного
  $n_f = -(d_1 k)_f = (d_1 A)_f/512 \in \mathbb{Z}$ образует дискретные топологические заряды материи.
\item \textbf{Единая обратимая траектория Вселенной $\symCapitalPsi^t = g^t(\mathrm{IV}(3))$:}
  поворот фазы есть сложение по модулю~$512$, $\varphi^{t+1} = (2\varphi^t - \varphi^{t-1} + \Phi^t)\bmod 512$
  (с буквальным обращением времени $\varphi^{t-1} = (2\varphi^t - \varphi^{t+1} + \Phi^t)\bmod 512$), стартующее из единственного
  $9$-битного начального вектора $\mathrm{IV}(3) = 2^{2^3} + (2^3-1) = 263$.
\end{enumerate}

Возникает главный вопрос: \emph{почему и каким образом из этих четырёх правил в целых числах~$\mathbb{Z}$ на макроуровне~$\Tlayer$
рождается вся классическая и современная теоретическая физика --- непрерывный математический анализ, уравнение Шрёдингера,
гидродинамика Навье--Стокса, неабелевы калибровочные поля Янга--Миллса $U(1)\times SU(2)\times SU(3)$ и общая теория относительности Эйнштейна?}
Ответ даёт следующая теорема: все перечисленные континуальные теории являются длинноволновыми проекциями тех же самых четырёх
правил \textbf{Z1--Z4} при прохождении целочисленного сигнала сквозь биномиальное окно макроприбора~$\mathcal{B}_N$ ($N\gg 1$).

\begin{theorem}[Вывод континуальных уравнений физики и щели массы Янга--Миллса из целочисленного автомата $g^t(\mathrm{IV}(3))$]
\label{thm:continuum-emergence-ym-mass-gap}
Пусть $\symCapitalPsi^t = g^t(\mathrm{IV}(3))$ --- целочисленная траектория на решётке кубооктаэдра~\latGCKshort\ ($d=3$),
а $\mathcal{B}_N$ --- оператор мягкого чтения прибора по окну из $N\gg 1$ шагов решётки с биномиальными весами путей
$w_N(k) = 2^{-2N}\binom{2N}{N+k}$ (\cref{def:soft-readout,thm:t-cr,thm:t-cl}). Тогда:
\begin{enumerate}[label=(\roman*)]
\item \textbf{Происхождение непрерывности, производных $\partial_\mu$ и констант $\pi,\ e,\ \planckHbar$:}
  при $N\gg 1$ дискретный сдвиг на один шаг решётки~$dh$ внутри биномиального окна $\mathcal{B}_N$ имеет относительный вес
  \[
    \frac{w_N(k+1) - w_N(k)}{w_N(k)} \;=\; -\frac{2k+1}{N+k+1} \;=\; -\frac{2x}{N\,dh^2}\,dh + O(N^{-1}),
    \qquad x := k\,dh,
  \]
  переводящий кограничный оператор $d_0$ в гладкий градиент $\nabla$, биномиальную огибающую $w_N(k)$ --- в гауссиан
  $(1/\sqrt{\pi N})\,e^{-k^2/N}$ (рождающий константы~$\pi$ и~$e$), кусочно-линейный $L^1$-ромб $|C(A)|+|S(A)|\equiv Q_\square$ ---
  в гладкую $L^2$-окружность $\cos^2\varphi + \sin^2\varphi = 1$, а полный оборот кольца $h = 512\,s_0$ при переходе к радианной мере $2\pi$ ---
  в приведённую постоянную Планка $\planckHbar = h/(2\pi)$. Непрерывность пространства-времени есть свойство апертуры прибора~$\mathcal{B}_N$, а не носителя~$\Mlayer$.
\item \textbf{Происхождение уравнений Шрёдингера и Навье--Стокса и отсутствие вихревого взрыва (глобальная регулярность):}
  применение $\mathcal{B}_N$ к ортогональному расщеплению голономии~\eqref{eq:holonomy-split} переводит нечётную часть~\eqref{eq:odd-holonomy-continuity}
  в макроскопическое уравнение непрерывности $\partial_t\rho + \operatorname{div}(\rho\mathbf{v}) = 0$, а чётный дефицит~\eqref{eq:even-holonomy-phase} ---
  в уравнение фазы Гамильтона--Якоби--Маделунга и уравнение импульса Навье--Стокса с кинематической вязкостью решётки $\nu_{\mathrm{CA}} = \caKappaGeom^2\,dh^2/dt$.
  При этом на носителе~$\Mlayer$ скорость переноса ограничена световым конусом $|\mathbf{v}|\le c_0 = dh/dt$, плотность ограничена
  потолком регистра $\mu\le K_P = 64$, а циркуляция скорости по минимальному замкнутому контуру грани ограничена порогом Паули $\Phi_{\mathrm{Pauli}}=256$ тиков,
  поэтому дискретная завихренность и энстрофия ограничены сверху всюду и при всех $t\in\mathbb{Z}$:
  \begin{equation}
  \label{eq:navier-stokes-regularity-bound}
    \bigl|\operatorname{curl}_{\mathrm{disc}}\mathbf{v}(x,t)\bigr| \;\le\; \frac{c_0}{dh} \;<\; \infty,
    \qquad
    \rho(x,t) \;=\; \frac{\mu(x,t)^2}{K_P^2} \;\le\; 1.
  \end{equation}
  Кажущаяся вихревая сингулярность континуальных уравнений Навье--Стокса при $dh\to 0$ является артефактом отбрасывания планковского шага~$dh$ и потолка $K_P=64$.
\item \textbf{Происхождение неабелевых полей Янга--Миллса $U(1)\times SU(2)\times SU(3)$, конфайнмента и строго положительной щели массы (Yang--Mills Mass Gap):}
  на клеточном комплексе одношаговой оболочки кубооктаэдра ($12$ вершин, $24$ ребра, $8\triangle\cup 6\square$ граней)
  остаток деления $A_{xy}$~\eqref{eq:quotient-remainder-512} определяет решёточную калибровочную связность, а кограница $B_f = (d_1 A)_f = 512\,n_f$ ---
  напряжённость калибровочного поля с тождеством Бьянки $d_2 B \equiv 0$~\eqref{eq:discrete-gauss-b}.
  Калибровочная группа Стандартной модели однозначно порождается тремя геометрическими компонентами оболочки кубооктаэдра в размерности $d=3$:
  \begin{equation}
  \label{eq:sm-gauge-group-from-cuboctahedron}
  \begin{aligned}
    \dim U(1) &\;=\; 1_{\mathrm{core}} \;=\; 1
      &&\text{(общая фаза звезды $N_{12}$, $\hat M = 1 + |N_{12}|\,|\triangle| = 97$)},\\
    \dim\mathfrak{su}(2) &\;=\; \frac{|\square|}{2} \;=\; d \;=\; 3
      &&\text{(три пары противоположных квадратных граней $\square$)},\\
    \dim\mathfrak{su}(3) &\;=\; |\triangle| \;=\; 2^d \;=\; 8
      &&\text{(восемь треугольных граней $\triangle$ --- глюонный октет)}.
  \end{aligned}
  \end{equation}
  В этом неабелевом секторе выполняются как \textbf{точный конфайнмент}, так и \textbf{строго положительная щель массы} $\Delta E_{\mathrm{YM}} > 0$:
  \begin{enumerate}[label=(\alph*)]
  \item по закону Гаусса на замкнутой оболочке кубооктаэдра~\eqref{eq:discrete-gauss-b} сумма топологических потоков по всем граням равна нулю,
    $\sum_{f\in 8\triangle\cup 6\square} n_f \equiv 0$, поэтому уединённое возбуждение одной треугольной грани ($n_\triangle = \pm 1$ при остальных $n_f = 0$,
    изолированный цветовой заряд или свободный глюон) топологически невозможно ($1\ne 0$): допустимы только нейтральные замкнутые конфигурации
    (дипольные мезонные пары $n_{\triangle_1}+n_{\triangle_2}=1-1=0$ и барионный трёхгранный угол $K_3$ на трёх гранях октанта);
  \item в силу порога Гейзенберга~\eqref{eq:heisenberg-divisibility-lock} любое подпороговое возмущение ворот $|\Phi| < \Delta\varphi_{\mathrm{disc}} = 41$ тик
    тождественно обращается в нуль ($\Phi \equiv 0$, $n_E = 0$), а минимальное ненулевое возбуждение калибровочного поля несёт ровно один квант действия $n_E = 1$
    ($41$ тик кольца) на микроуровне и минимальный замкнутый плакетный скачок $|n_\triangle|\ge 1$ ($512$ тиков) в неабелевом канале треугольных граней,
    что на макроуровне~$\Tlayer$ задаёт строго положительную щель массы адронного сектора (натяжение цветовой струны внутри~$K_3$, \cref{prop:ckm-pmns-sectors}):
    \begin{equation}
    \label{eq:yang-mills-mass-gap}
      \Delta E_{\mathrm{YM}}^{(\Mlayer)} \;\ge\; \ladderEnergy \;=\; \frac{\Delta\varphi_{\mathrm{disc}}}{\htick} \;>\; 0,
      \qquad
      \Delta m_{\mathrm{YM}}^{(\Tlayer)} \;\ge\; \Lambda_{\mathrm{QCD}} \;=\; \frac{m_p}{d_{\mathrm{in}}(\square)} \;\approx\; 234{,}6\ \mathrm{МэВ} \;>\; 0.
    \end{equation}
  \end{enumerate}
\item \textbf{Происхождение метрики Эйнштейна $g_{\mu\nu}$ и гравитации без сингулярностей:}
  вес ячейки $q(x)=\mu^2(x)$ входит в знаменатель ворот $D(x) = \zeta_{\mathrm{real}}(x) + \mu^2(x) + K_P$, замедляя локальный темп хода фазовых часов
  в $n(x) = D(x)/D_{\mathrm{vac}} = 1 + \varepsilon_{\mathrm{eff}}(x) \ge 1$ раз, а одношаговое реберное преломление $\Delta_t A_\perp = (d_0\Phi)_\perp < 0$~\eqref{eq:local-edge-refraction}
  при сглаживании $\mathcal{B}_N$ суммируется в эффективный метрический интервал Общей теории относительности
  \[
    ds^2 \;=\; -g_{00}(x)\,c^2 dt^2 + g_{ij}(x)\,dx^i dx^j,
    \qquad
    g_{00}(x) \;=\; \frac{1}{n(x)^2} \;=\; \left(\frac{D_{\mathrm{vac}}}{D(x)}\right)^{\!2} \;\in\; \left[\frac{D_{\mathrm{vac}}^2}{D_{\max}^2},\; 1\right],
  \]
  где конечный потолок регистра $\mu\le K_P = 64$ ограничивает знаменатель величиной $D(x) \le D_{\max} = (|N_{12}|+1)K_P^2 + K_P = 13\times 4096 + 64 = 53312 < \infty$,
  тождественно устраняя гравитационные сингулярности $g_{00}\to 0$ ($D\to\infty$).
\end{enumerate}
\end{theorem}

\begin{proof}
Пункт~\emph{(i)} следует из разложения отношения биномиальных коэффициентов $\binom{2N}{N+k+1}/\binom{2N}{N+k} = (N-k)/(N+k+1)$ и теоремы~\ref{thm:t-cl}.
В пункте~\emph{(ii)} подстановка $z = \mu e^{i\varphi}$ в~\eqref{eq:holonomy-split} после усреднения~$\mathcal{B}_N$ даёт уравнения Маделунга--Навье--Стокса (\cref{thm:t-cr,lem:bohm-chi}), а оценки~\eqref{eq:navier-stokes-regularity-bound} следуют из $c_0 = dh/dt$ (\Axref{ax:A1}) и $\mu\in\{1,\dots,64\}$ (\Axref{ax:A7}).
В пункте~\emph{(iii)} размерности~\eqref{eq:sm-gauge-group-from-cuboctahedron} фиксированы числами элементов кубооктаэдра~\eqref{eq:cubocta-face-invariants} при $d=3$: $|\square|/2 = d = 3 = \dim\mathfrak{su}(2)$ и $|\triangle| = 2^d = 8 = \dim\mathfrak{su}(3)$; невозможность изолированного заряда $n_\triangle = \pm 1$ следует из $\sum_f n_f \equiv 0$~\eqref{eq:discrete-gauss-b}, а неравенства~\eqref{eq:yang-mills-mass-gap} --- из затвора $|\Phi|\ge 41$ и целочисленности $n_\triangle\in\mathbb{Z}$.
В пункте~\emph{(iv)} монотонность $D(x) \ge D_{\mathrm{vac}} = 82$ и верхняя грань $D_{\max} = 13\cdot 64^2 + 64 = 53312$ при $\mu\le 64$ дают строго положительную нижнюю границу $g_{00}(x) \ge (82/53312)^2 > 0$ в каждой ячейке решётки.
\end{proof}

\begin{table}[htbp]
\centering
\caption{Как из четырёх правил в целых числах~$\mathbb{Z}$ на~$\Mlayer$ при сглаживании прибором~$\mathcal{B}_N$ строится вся континуальная теоретическая физика на~$\Tlayer$.}
\label{tab:z-to-continuum-physics}
\footnotesize
\setlength{\tabcolsep}{3.5pt}
\begin{tabular}{@{}p{0.23\linewidth}p{0.26\linewidth}p{0.25\linewidth}p{0.20\linewidth}@{}}
\toprule
Правило в целых числах~$\mathbb{Z}$ на~$\Mlayer$ & Что происходит при сглаживании~$\mathcal{B}_N$ & Теория на~$\Tlayer$ (учебник) & Какая проблема континуума снимается \\
\midrule
Замок $1+2^d=d^2\implies d=3$, кубооктаэдр $|N_{12}|=12$, $8\triangle+6\square$, регистр $9$ бит ($\mathbb{Z}_{512}$) &
Биномиальный подсчёт путей $\binom{2N}{N+k}/2^{2N}\to$ гауссиан, $d_0\to\partial_\mu$, $|C|+|S|=128\to\cos^2+\sin^2=1$ &
Непрерывное $\mathbb{R}^{3,1}$, матанализ, константы $\pi,\ e,\ \planckHbar=h/(2\pi)$ &
Произвол размерности $d=3$ и аксиомы континуума \\
\addlinespace[2pt]
Расщепление голономии: чётный дефицит $\delta\zeta_{\mathrm{real}}\ge 0$ и антисимметричный поток $\zeta_{\mathrm{imag}}=-\operatorname{div}_{\mathrm{disc}}\mathbf{j}$ &
$\zeta_{\mathrm{imag}}\to\partial_t\rho+\operatorname{div}(\rho\mathbf{v})=0$; $\delta\zeta_{\mathrm{real}}\to -\frac{\planckHbar^2}{2m}\nabla^2\psi$ и давление Бома при $\nu_{\mathrm{CA}}=\caKappaGeom^2 dh^2/dt$ &
Уравнение Шрёдингера, гидродинамика Маделунга и Навье--Стокса &
Вихревой взрыв Навье--Стокса ($|\operatorname{curl}\mathbf{v}|\le c_0/dh$, $\mu\le 64$) \\
\addlinespace[2pt]
Деление с остатком $(d_0\varphi)_{xy}=512\,k_{xy}+A_{xy}$ на гранях $8\triangle\cup 6\square$, $B_f=(d_1 A)_f=512\,n_f$, $|\Phi|\ge 41$ &
Остаток $A_{xy}\to$ потенциал $A_\mu$; частное $n_f=-(d_1 k)_f\to$ ток источника и поля $U(1)\times SU(2)_{3\square}\times SU(3)_{8\triangle}$ &
Уравнения Максвелла и калибровочная теория Янга--Миллса, КЭД и КХД &
Внешние источники, УФ-расходимости, конфайнмент ($\sum n_f=0$) и щель массы ($\Delta m_{\mathrm{YM}}>0$) \\
\addlinespace[2pt]
Вес $q=\mu^2$ в знаменателе ворот $D(x)=\zeta_{\mathrm{real}}+\mu^2+64$ и реберный поворот $\Delta_t A_\perp=(d_0\Phi)_\perp < 0$ &
Локальное замедление такта $n(x)=D(x)/D_{\mathrm{vac}}\to g_{00}=n^{-2}$, сумма реберных поворотов $\to$ принцип Ферма и геодезические &
Общая теория относительности Эйнштейна, преломление света $\theta=4GM/(c^2 b)$ &
Гравитационные сингулярности ($D\le 53312<\infty$) и несовместимость ОТО с КМ \\
\bottomrule
\end{tabular}
\end{table}

\section{Что было закрыто ревизией единиц и что остаётся на границе \texorpdfstring{$\Mlayer/\Tlayer$}{M/T}}

Перечень проверок, отделивших свойства носителя от ошибок записи и перевода единиц.

\begin{enumerate}[leftmargin=*]
\item \textbf{Откуда брался кажущийся источник в декартовой «лягушке» и рост до $\rho\approx 63$.}
  В прежнем замере при поднятой амплитуде сида суммарная $\sum_x|z|^2$ росла
  с $768$ до $2{,}3\times10^{4}$ за $512$ тактов и упиралась в $\rho\approx 63$ на ячейку по двум причинам:
  во-первых, линейная комбинация $2Z(t)-Z(t-1)+\delta Z$ в декартовых дорожках $(U,V)$ не сохраняет
  радиус $\mu$; во-вторых, в затворе потолка Бекенштейна (\texttt{bekenstein\_scale\_spinor})
  стоял лишний множитель $2^{\mathrm{frac\_bits}}=64$ (\texttt{1 << (2 * fb)} вместо \texttt{1 << fb}),
  поднимавший потолок~\Axref{ax:A7} с $\rho_{\max}=1$ до $\rho_{\max}=64$.
  После исправления потолка на $\rho_{\max}=1$ ($\mu_{\max}=64$ кванта) и перехода к расщеплению
  голономии~\eqref{eq:holonomy-split} источник исчезает, а на первой координационной сфере $S_1$
  канонического 3D-вихря ровно $100\%$ из $12$ узлов кубооктаэдра лежат над порогом Гейзенберга
  ($\max\Delta\varphi = 0{,}833\ \text{рад} \ge 0{,}5\ \text{рад}$, на $S_3$ --- $0\%$, $\max\Delta\varphi = 0{,}382\ \text{рад}$),
  воспроизводя радиус одевания $R_{\mathrm{dress}}=1\cdot dl$.
\item \textbf{Граница единиц внутри ворот $\Phi$: сдвиг \texttt{>> frac\_bits} и расхождение шкал $32$ против $41$ (\cref{tab:m-t-bridge}).}
  В формуле ворот $\Phi$ отношение голономии к знаменателю $\zeta_{\mathrm{real}} + R^2 + K_P$
  определяет угол поворота в радианах только тогда, когда числитель и знаменатель выражены
  в одних и тех же единицах квадрата квантов ($\mathrm{quanta}^2$, без предварительного усечения \texttt{>> frac\_bits}).
  При этом умножение безразмерного угла $\theta\ [\text{рад}]$ на $K_P = 2^{\mathrm{frac\_bits}} = 64$
  переводит радианы в дробные доли фиксированной точки Q3.6 ($1\ \text{рад} = 64$, где порог Гейзенберга
  $\Delta\varphi_{\min}=1/2\ \text{рад}$ равен $32$), тогда как таблица поворотов и порог
  $\Delta\varphi_{\mathrm{disc}} = 41$ живут в тиках кольца $\mathbb{Z}_{512}$ ($1\ \text{рад} = 512/(2\pi)\approx 81{,}5$ тика,
  где тот же угол $1/2\ \text{рад}$ равен $41$ тику). Именно эта несоизмеримость двух шкал ($32$ в Q3.6 против $41$ в $\mathbb{Z}_{512}$)
  вместе с потерей чётного дефицита $\delta\zeta_{\mathrm{real}}$~\eqref{eq:even-holonomy-phase} и второй компоненты спинора
  создавала иллюзию, будто гладкие волны и вихри «насыщаются на $32$ и не достают до $41$».
\item \textbf{Условие~\Axref{ax:A3} в глобальной форме не выводится и не выполняется.}
  Оно сохраняется в тексте как историческая запись и подлежит переформулировке:
  «информация не теряется» (теорема~\ref{thm:info-bijection}) плюс
  локальное уравнение непрерывности (предложение~\ref{prop:source-or-flux}).
\item \textbf{Точность $\mathbb{Z}_4$ — только в декартовой записи.}
  Арифметический сдвиг вправо есть floor, а floor нечётен; при несимметричном
  округлении поворот теряет эквивариантность (порядка $0{,}15\%$ дорожек за такт).
  Симметричное округление возвращает точное равенство
  $\mathrm{step}(\hat g f)=\hat g\,\mathrm{step}(f)$ для $\hat g\in\mathbb{Z}_4$.
  В полярном регистре вопроса нет вовсе: поворот там --- сложение по модулю,
  и он точен для всех $\phaseRingCount$ углов без округления.
\end{enumerate}

\begin{deadend}{порог ворот: оценка «порядка единицы» против замера}
Оценка по формуле давала $\Phi$ порядка единицы при $\rho\le1$; отсюда был сделан вывод,
что порог $\Delta\varphi_{\mathrm{disc}}=41$ недостижим «в сорок раз».
Прямой замер поля $\Phi$ на каноническом сиде дал сырой максимум $22$ тика ---
недобор меньше чем вдвое, а не в сорок раз. Оценка была неверна арифметически;
существо вопроса при этом осталось открытым.
След: \texttt{projected\_phi\_int} на сиде $\mathrm{IV}=263$, $24^3$, до затвора
и после; при амплитуде $\times1{,}5$ срабатывает одна ячейка, при $\times2{,}45$ ---
шестьдесят семь, и $\Phi$ насыщается на $256$.
\end{deadend}

\section{Демаркация и что опровергнет эту главу}

\begin{enumerate}[leftmargin=*]
\item \textbf{Точность $\mathbb{Z}_4$:} побитовая проверка
  $\mathrm{step}(\hat g f)=\hat g\,\mathrm{step}(f)$ для $\hat g\in\{\pm1,\pm i\}$
  на любом состоянии. Ненулевое расхождение при точной целочисленной арифметике
  опровергает предложение~\ref{prop:z4-exact}.
\item \textbf{Природа дрейфа $q$:} дрейф должен убывать при росте разрешения таблицы
  поворотов (удвоение разрядности LUT). Если он от разрешения не зависит ---
  источник не в таблице, и предложение~\ref{prop:q-scope} неверно.
\item \textbf{Источник в непрерывности:} поячеечный остаток
  $\Delta q+\operatorname{div}j$. Если он нулевой при любой схеме ---
  величина сохраняется локально, и пункт~1 перечня «не выведено» снимается.
\end{enumerate}

\noindent
Итог главы в одну строку: \emph{носитель сохраняет не величину, а различимость},
и всякое «сохранение» на~$\Tlayer$ есть длинноволновый предел этой различимости
при конечном разрешении таблицы.
